% ch24: 积分法
\chapter{积分法——换元、分部、有理函数}
\begin{applicationbox}
学完本章：掌握换元积分法、分部积分法、有理函数积分
\end{applicationbox}

\section{换元积分法} \(\int f(g(x))g'(x)dx = \int f(u)du\)（令 \(u=g(x)\)）

\begin{examplebox}[1] \(\int 2x\cos(x^2)dx = \sin(x^2)+C\)（令 \(u=x^2\)）\end{examplebox}

\section{分部积分法} \(\int u dv = uv - \int v du\)

\begin{examplebox}[2] \(\int x e^x dx = x e^x - e^x + C\)\end{examplebox}

\section{有理函数积分} 分解为部分分式后逐项积分。

\section{习题}
\begin{exercise}[0] 换元法的公式是什么？\end{exercise}
\begin{exercise}[0] 分部积分法的公式是什么？\end{exercise}
\begin{exercise}[0] \(\int e^x dx = ?\)\end{exercise}
\begin{exercise}[0] \(\int \frac1x dx = ?\)\end{exercise}
\begin{exercise}[1] \(\int (2x+1)^3 dx\)\end{exercise}
\begin{exercise}[1] \(\int \cos 3x dx\)\end{exercise}
\begin{exercise}[1] \(\int xe^{x^2} dx\)\end{exercise}
\begin{exercise}[1] \(\int x\cos x dx\)\end{exercise}
\begin{exercise}[2] \(\int \sin^2 x \cos x dx\)\end{exercise}
\begin{exercise}[2] \(\int \ln x dx\)\end{exercise}
\begin{exercise}[2] \(\int \frac{x}{x^2+1} dx\)\end{exercise}
\begin{exercise}[2] \(\int \arctan x dx\)\end{exercise}
\begin{exercise}[3] \(\int \frac{\ln x}{x} dx\)\end{exercise}
\begin{exercise}[3] \(\int e^x \sin x dx\)\end{exercise}
\begin{exercise}[3] \(\int \frac{dx}{x^2-1}\)\end{exercise}
\begin{exercise}[3] \(\int x^2 e^x dx\)\end{exercise}
\begin{exercise}[3] \(\int \tan x dx\)\end{exercise}
\begin{exercise}[4] \(\int \sqrt{1-x^2} dx\)\end{exercise}
\begin{exercise}[4] \(\int \sin^3 x dx\)\end{exercise}
\begin{exercise}[4] \(\int \frac{dx}{x^2+2x+2}\)\end{exercise}
\begin{exercise}[4] \(\int e^{2x}\cos 3x dx\)\end{exercise}
\begin{exercise}[4] \(\int \frac{x}{(x+1)(x+2)} dx\)\end{exercise}
\begin{exercise}[5] \(\int \sec^3 x dx\)\end{exercise}
\begin{exercise}[5] \(\int \frac{dx}{(x^2+1)^2}\)\end{exercise}
\begin{exercise}[5] \(\int \frac{x^2}{(x^2+1)^2} dx\)\end{exercise}
\begin{exercise}[5] \(\int \sin(\ln x) dx\)\end{exercise}
\begin{exercise}[6] \(\int \frac{x^4}{x^4-1} dx\)\end{exercise}
\begin{exercise}[6] \(\int \frac{dx}{x\sqrt{x^2-1}}\)\end{exercise}
\begin{exercise}[6] \(\int e^{ax}\sin bx dx\) 的递推公式。\end{exercise}
\begin{exercise}[7] 建立 \(\int_0^{\pi/2} \sin^n x dx\) 的递推公式。\end{exercise}
